% HD-N21. GPT-6 Astra Pro /root/imprimitive, 21 November 2026.
% Inherited effort; numerical setting not exposed.
\section{HD-N21: abelian-support sums and invariant trace kernels}
\label{sec:abelian-support-weil-obstruction}

Let $(A,K)$ be a complex abelian variety of Weil type of dimension
$2n$ and signature $(n,n)$, where $n\ge2$. Let
\[
 j_\nu:B_\nu+a_\nu\hookrightarrow A
\]
be actual inclusions of translates of abelian subvarieties of
dimension $n$. On each $B_\nu+a_\nu$ take an actual coherent
sheaf $E_\nu$ of positive generic rank, and take integers $b_\nu$.
Consider the actual perfect complex
\[
          F=\bigoplus_{\nu=1}^r j_{\nu,*}E_\nu[b_\nu].
                                                               \tag{N21.1}
\]
Every such complex whose whole character is a nonzero rational
Weil class
\[
                   \operatorname{ch}(F)=t\ne0,\qquad
                             t\in W_K(A)                      \tag{N21.2}
\]
has a degree-two categorical trace kernel of dimension at least
$n-1$. The same lower bound holds on invariant Ext for every
actual coherent action by a faithful finite point subgroup
$G\subset A\times\widehat A$. Hence none of these objects is
weakly $G$-semiregular.

All multiplicities, shifts, intersections of supports, and mixing
linearizations are allowed in this statement. No assumption of
negative self-Ext vanishing is used: the asserted trace kernel
exists also for members of this class which are gluable.
In particular it applies directly to signed sums of translated
graph structure sheaves on a CM product. It does not replace a
complex having new nonzero differentials by such a direct sum.

\paragraph{1. The mixed forms and the exact trace map.}
Write
\[
 T_+=H^{0,1}(A)_\sigma,\qquad
 T_-=H^{0,1}(A)_{\bar\sigma},\qquad
                         \mathcal M=T_+\wedge T_- .
                                                               \tag{N21.3}
\]
Both $T_\pm$ have dimension $n$, and $\mathcal M$ is the
$n^2$-dimensional tensor summand of $H^2(A,\mathcal O_A)$.
The two complex Weil lines are
\[
 \bigwedge^nH^{1,0}(A)_\sigma\otimes\bigwedge^nT_+,\qquad
 \bigwedge^nH^{1,0}(A)_{\bar\sigma}\otimes\bigwedge^nT_-.
                                                               \tag{N21.4}
\]
Consequently $\eta t=0$ for every $\eta\in\mathcal M$:
one of its two antiholomorphic factors repeats the full
antiholomorphic volume in each line.

The actual scalar characteristic morphism is
\[
 \mathcal M\longrightarrow\operatorname{Ext}^2(F,F),\qquad
                       \eta\longmapsto\eta\operatorname{id}_F.
                                                               \tag{N21.5}
\]
Naturality and the Atiyah--Chern trace identity give
\[
 \sigma_F^2(\eta\operatorname{id}_F)
                  =\eta\,\operatorname{ch}(F)=0.              \tag{N21.6}
\]
Here all components of the categorical trace are included.
The precise published characteristic and trace identities are
Buchweitz--Flenner, Adv. Math.217 (2008),243--281,
Proposition6.2.1, Corollary6.3.2 and Theorem6.4.2,
author version \href{https://arxiv.org/pdf/math/0606730}
{arXiv:math/0606730}, pp.29--31, at the primary locators
already checked in HD-N19.

Every point autoequivalence $T_a^*(-)\otimes P$ fixes
$H^2(A,\mathcal O_A)$. For translation this follows from
invariant Dolbeault forms; for flat tensoring it follows
from its scalar action and its zero Atiyah class.
The characteristic transformation is natural. Thus every
class (N21.5) is invariant under every coherent linearization
on $F$, even when the linearization mixes the displayed
summands. No invariant-projector choice is required.

\paragraph{2. The actual complement isogeny detects the scalar classes.}
Fix one support $B+a$ and write $j$ for its inclusion.
The restriction
\[
             j^*:H^{0,1}(A)\longrightarrow H^{0,1}(B)
                                                               \tag{N21.7}
\]
is surjective. One exact way to see all maps at once is
Poincare reducibility: for a fixed polarization $L$ on $A$
and the origin-preserving inclusion $i:B\hookrightarrow A$,
put
\[
 C=\ker(i^\vee\lambda_L:A\to\widehat B)^0,\qquad
 p:B\times C\longrightarrow A,\quad p(x,y)=x+y.
                                                               \tag{N21.8}
\]
This is an actual isogeny. The precise primary author-text
locator, including this kernel and the addition map, is
Milne, \emph{Abelian Varieties}, version2.00, ChapterI,
Proposition10.1 and its proof, printed p.42:
\url{https://www.jmilne.org/math/CourseNotes/AV.pdf}.
The complete stated proof was read. Pullback on $H^1$
is an isomorphism, and restriction to the first factor
proves (N21.7).

The inverse image of $B+a$ under $p$ is a disjoint union
of translates $B\times\{c\}$. On each component the
restriction of $p$ to $B$ is a translation isomorphism.
Finite flat base change therefore identifies the actual
pullback $p^*j_*E$ with the direct sum of the corresponding
sheaves $E'\boxtimes\mathcal O_c$. For each component the
derived Kunneth morphism gives
\[
 \operatorname{Ext}^2_B(E',E')\otimes
          \operatorname{Hom}_C(\mathcal O_c,\mathcal O_c)
       \lhook\joinrel\longrightarrow
 \operatorname{Ext}^2_{B\times C}
       (E'\boxtimes\mathcal O_c,E'\boxtimes\mathcal O_c).
                                                               \tag{N21.9}
\]
It is the indicated direct summand of the actual tensor
product of the two derived Hom complexes.
The $(2,0)$ Kunneth component of the pullback of
$\eta\operatorname{id}_{j_*E}$ in (N21.9) is exactly
\[
                    (j^*\eta)\operatorname{id}_{E'}
                                  \otimes\operatorname{id}_{\mathcal O_c}.
                                                               \tag{N21.10}
\]
Its trace on $B$ has lowest component
$\operatorname{rank}(E)\,j^*\eta$. Thus nonzero $j^*\eta$
makes (N21.10) nonzero, also when $E$ is not locally free:
it is perfect on the smooth $B$, and the rank term of
its Atiyah--Chern trace is still its positive generic rank.

For completeness, finite-etale pullback on Ext is faithful.
Under its actual adjunction it is induced by
$\mathcal O_A\to p_*\mathcal O_{B\times C}$; the trace
divided by $\deg p$ splits this unit, so the induced map
on each Ext group is injective. Naturality also sends a
zero scalar class to zero after pullback. Hence
the kernel of the scalar map lies in the kernel of every
restriction $j_\nu^*$. The reverse assertion also has an
actual map: derived pullback sends the scalar morphism
$\mathcal O_A\to\mathcal O_A[2]$ to its ordinary coherent
restriction on $B$, and the projection formula identifies
its action on $j_*E$ with the pushforward of that restricted
scalar action on $E$. Thus restriction zero makes this
diagonal scalar action zero. We have the exact equality
\[
 \ker\bigl(\mathcal M\xrightarrow{\eta\operatorname{id}}
                       \operatorname{Ext}^2(F,F)\bigr)
     =\bigcap_\nu
       \ker\bigl(\mathcal M\xrightarrow{j_\nu^*}
                     H^2(B_\nu,\mathcal O_{B_\nu})\bigr).
                                                               \tag{N21.11}
\]
The direct-sum blocks are actual blocks, and shifts cause
no sign in degree two.

\paragraph{3. A mixed support contributes at least $n-1$ dimensions.}
Put $U=j^*T_+$ and $V=j^*T_-$ inside $H^{0,1}(B)$.
Equation(N21.7) gives $U+V=H^{0,1}(B)$. The image of
$j^*\mathcal M$ is exactly $U\wedge V$.
If $U,V$ are both nonzero, put
$a=\dim U$, $b=\dim V$, $c=\dim(U\cap V)$, so
$a+b-c=n$. Choosing a basis of the intersection and
complementary bases gives the exact direct-sum wedge map
and dimension
\[
                 \dim(U\wedge V)=ab-\frac{c(c+1)}2
                                               \ge n-1.     \tag{N21.12}
\]
Here is the elementary inequality with all endpoints retained.
For $c=0$ its difference from $n-1$ is
$(a-1)(b-1)\ge0$. For $c\ge1$, write $u=a-c$ and
$v=b-c$. The difference is
\[
          uv+(c-1)(u+v)+\frac{(c-1)(c-2)}2\ge0.
\]
It follows from (N21.11) that one such support gives
at least $n-1$ independent classes in (N21.5), and every
one of them belongs to the invariant trace kernel.

\paragraph{4. The alternative is excluded by the exact exterior Gysin map.}
It remains to prove that at least one mixed support occurs.
Suppose instead that $j^*\mathcal M=0$ for every support.
If both $U$ and $V$ were nonzero, choosing a nonzero
element of either and using $U\wedge V=0$ would put both
spaces in the same line. Their sum could not have dimension
$n\ge2$. Therefore, on each support, one of $T_+,T_-$
restricts to zero and the other maps isomorphically onto
$H^{0,1}(B)$.

The following describes the actual pushforward, rather than
inferring a statement from an unspecified stability property.
For definiteness suppose $T_+$ restricts to zero.
Complex conjugation gives
\[
 \begin{split}
 \ker\bigl(H^1(A,\mathbf C)\to H^1(B,\mathbf C)\bigr)
        &=H^{1,0}(A)_{\bar\sigma}\oplus T_+,\\
 H^{1,0}(A)_\sigma\oplus T_-&\ \xrightarrow{\ \sim\ }\ H^1(B,\mathbf C).
 \end{split}                                                  \tag{N21.13}
\]
Let $q:A\to A/B$ be the actual quotient. An oriented
integral basis $\alpha_1,\ldots,\alpha_{2n}$ of
$H^1(A/B,\mathbf Z)$, normalized by
$\int_{A/B}\alpha_1\cdots\alpha_{2n}=1$, gives
\[
 [B]=q^*(\alpha_1\cdots\alpha_{2n}).
                                                               \tag{N21.14}
\]
Indeed $B$ is the fibre over a point of this smooth quotient,
so this is its ordinary Gysin fundamental class. In
(N21.13) it is the normal volume, with $n$ factors from
each of the two displayed kernel summands.

Give the $\sigma$ and $\bar\sigma$ eigenspaces circle weights
$+1$ and $-1$. Equation(N21.14) has weight zero. Inverting
the second map in (N21.13) gives an actual exterior-algebra
section $s$ of $j^*$ on complex cohomology. A $(k,k)$ class
on $B$ lifts under $s$ to a class with $k$ factors of weight
$+1$ and $k$ factors of weight $-1$, hence of weight zero.
The projection formula is the exact exterior Gysin formula
\[
                  j_*(\gamma)=[B]\wedge s(\gamma).
                                                               \tag{N21.15}
\]
Both tangent Todd classes are one. GRR, at the exact Fulton
Theorem15.2 locator in HD-N20, for this actual
embedding therefore gives
$\operatorname{ch}(j_*E)=j_*\operatorname{ch}(E)$, and
(N21.15) shows that every homogeneous component has
weight zero. The other choice of vanishing eigenspace
exchanges the two signs and gives the same result.
Translations and shifts do not alter this conclusion.

Thus the entire character of (N21.1) would have weight
zero. But the two nonzero-weight spaces (N21.4) have
weights $2n$ and $-2n$. Their sum has zero intersection
with the weight-zero space. This contradicts (N21.2).
Consequently a mixed support exists. Paragraph3 proves
the claimed lower bound.

\paragraph{5. Weak semiregularity, actual line twists and scope.}
The ambient trace kernel just constructed is invariant
under every coherent point action, by paragraph1. For a
faithful finite point subgroup the actual forgetful
Hochschild map is an isomorphism: the nonidentity
translation graph misses the diagonal, and a nontrivial
flat-tensor sector has no cohomology. Induction followed
by forgetting is $|G|$ on the identity sector. These are
the complete maps proved in HD-N12, whose proof works
in dimension $2n$. They identify weak $G$-semiregularity
with injectivity of the ambient trace on invariant Ext.
Hence the constructed invariant kernel proves the
stated weak-semiregularity failure.

There is also a literal common-line-twist extension.
For an actual line bundle $L$ on $A$, tensoring is an
equivalence and carries (N21.5) to the same scalar class
for $F\otimes L$. Its trace is multiplied by
$e^{c_1(L)}$, an invertible exterior class. An arbitrary
point action on $F\otimes L$ is conjugated to one on
$F$ by this equivalence; the actual point map is
\[
 (a,P)\longmapsto
       \bigl(a,P\otimes T_a^*L\otimes L^{-1}\bigr)
                                                               \tag{N21.16}
\]
when conjugating by tensoring with $L^{-1}$.
It is an automorphism of $A\times\widehat A$ and retains
finite faithfulness and the given coherence. Thus the
same obstruction holds for these actual common twists,
whose character is $e^{c_1(L)}t$. No fractional or
arbitrary real twist is replaced by a geometric functor.

For sixfolds the universal lower bound is two, and for
eightfolds it is three. This closes the indicated direct
abelian-support construction, including its actual
common line twists and every coherent faithful point
action. It does not exclude new extensions, cones or
other objects with the same virtual cycle expression.

\paragraph{6. A concrete gluable CM graph source and its exact scalar kernel.}
The preceding source class has explicit members with the
required whole character. Take
$A=E_i^n\times E_i^n$, where $E_i$ is the Gaussian CM
elliptic curve, and fix any positive integer $d$.
The actual endomorphism and polarization
\[
 J(P,Q)=(-Q,dP),\qquad J^2=-d,\qquad
                       \theta=dH_X+H_Y
                                                               \tag{N21.17}
\]
give a balanced action of $K=\mathbf Q(s)$, $s^2=-d$.
On each pair of holomorphic tangent lines the eigenvalues
are $s,-s$, so the signature is $(n,n)$.
The matrix identity $J^t\operatorname{diag}(d,1)
=-\operatorname{diag}(d,1)J$ verifies the required
Rosati relation for the actual polarization.

Choose integral degree-one forms $x_i,y_i$ on the first
elliptic factors and $u_i,v_i$ on the second, with
$x_i+\sqrt{-1}y_i$ and $u_i+\sqrt{-1}v_i$ holomorphic
for $E_i=\mathbf C/(\mathbf Z+\sqrt{-1}\mathbf Z)$,
and $x_iy_i,u_iv_i$ the positive point classes.
For $\epsilon\in\{-1,0,1\}^n$, let
$\Gamma_\epsilon$ be the actual graph
$Q_i=[\epsilon_i]P_i$. The quotient
$(P,Q)\mapsto(Q_i-[\epsilon_i]P_i)_i$ has this connected
kernel, hence
\[
 [\Gamma_\epsilon]
       =\prod_{i=1}^n g_{\epsilon_i,i},\qquad
       g_{t,i}=(u_i-tx_i)(v_i-ty_i).
                                                               \tag{N21.18}
\]
Both tangent bundles have Todd class one, so
$\operatorname{ch}(\mathcal O_{\Gamma_\epsilon})
=[\Gamma_\epsilon]$ by the actual embedding GRR map.

Put
\[
 \Omega_{s,i}=(u_i-sx_i)(v_i-sy_i),\qquad
 b_0=2(d+1),\quad b_+=s-d,\quad b_-=-s-d .
\]
Direct expansion gives the exact interpolation identity
\[
               2\Omega_{s,i}
                    =b_0g_{0,i}+b_+g_{1,i}+b_-g_{-1,i}.
                                                               \tag{N21.19}
\]
Define integers
$c_\epsilon=\operatorname{Re}\prod_i b_{\epsilon_i}$:
the even powers of $s$ are integer powers of $-d$ and
the odd powers have zero real part. Equations(N21.18)--(N21.19)
give
\[
 \sum_\epsilon c_\epsilon[\Gamma_\epsilon]
       =2^{n-1}\left(\prod_i\Omega_{s,i}
                                  +\prod_i\Omega_{-s,i}\right)
                              =:t_d\ne0.                     \tag{N21.20}
\]
The pullback action of $J$ sends $x_i$ to $-u_i$ and
$u_i$ to $dx_i$, and likewise sends $y_i$ to $-v_i$
and $v_i$ to $dy_i$. Thus the two one-forms in
$\Omega_{s,i}$ have eigenvalue $s$. The two products in
(N21.20) are exactly the conjugate top $K$-eigenvolumes,
so $t_d$ is a rational integral Weil class for the actual
action (N21.17). Its coefficient of $\prod_i u_iv_i$
is $2^n$, proving it is nonzero.

The actual source object is
\[
 F_d=
 \bigoplus_{c_\epsilon>0}
       \mathcal O_{\Gamma_\epsilon}^{\oplus c_\epsilon}
 \ \oplus\
 \bigoplus_{c_\epsilon<0}
       \mathcal O_{\Gamma_\epsilon}[1]^{\oplus(-c_\epsilon)}.
                                                               \tag{N21.21}
\]
Its whole character is (N21.20). It is gluable:
ordinary negative Ext between sheaves is zero; the only
additional possible negative group between the two shift
levels is a Hom between distinct supports. This Hom is
zero, since the target structure sheaf is torsion-free
on its integral support and the ideal of a distinct
equal-dimensional support acts there by a nonzero ideal.
Opposite coefficient signs never occur on the same
graph, by the definition of $c_\epsilon$.

This actual example has a larger exact scalar-kernel image.
Put $a_i=x_i-\sqrt{-1}y_i$ and $b_i=u_i-\sqrt{-1}v_i$,
the antiholomorphic forms of the CM elliptic factors.
The bases of $T_\pm$ are
$z_i^+=b_i-sa_i$ and $z_i^-=b_i+sa_i$.
On $\Gamma_\epsilon$ their restrictions are
$(\epsilon_i-s)a_i$ and $(\epsilon_i+s)a_i$.
Each diagonal mixed form $z_i^+\wedge z_i^-$ restricts
to zero on every graph. Conversely the graph $\Gamma_0$
occurs with coefficient $[2(d+1)]^n$, and each
$\Gamma_{e_i}$ occurs with the nonzero coefficient
$-d[2(d+1)]^{n-1}$. For $i<j$, restriction of
$z_i^+\wedge z_j^-$ and $z_j^+\wedge z_i^-$ to these
two graphs has the coefficient matrix
\[
              \begin{pmatrix}d&-d\\d+s&-d+s\end{pmatrix},
                    \qquad \det=2ds\ne0.                  \tag{N21.22}
\]
The different pairs $i<j$ give independent exterior
basis coefficients. Hence the joint restriction has
kernel precisely the $n$ diagonal forms, and
(N21.11) proves
\[
 \dim\operatorname{im}
       \bigl(\mathcal M\to\operatorname{Ext}^2(F_d,F_d)\bigr)
                                =n(n-1).                    \tag{N21.23}
\]
Every one of these classes is an invariant trace-kernel
class for every coherent faithful point action on $F_d$.
In particular the actual gluable sixfold object has at
least six such classes, and the eightfold object has at
least twelve. This supplies the exact failed trace map
for the exhibited CM construction.
